% !TeX root = ../概率论与数理统计笔记.tex
% 由主文档通过 \input 载入；请编译主文档。
\setcounter{chapter}{8} % 保留原书章号。
\chapter{概率不等式、大数定律与中心极限定理}
\label{chap:ps-limits}

本章解释样本平均为什么能够反映总体，以及误差为何常用正态分布近似。
先修内容为期望、方差与随机变量序列；学习目标是区分误差界、收敛类型和极限分布，并逐项核对定理条件。

\section{基本概率不等式}
\label{sec:ps-probability-inequalities}
证明马尔可夫不等式（Markov's inequality）与切比雪夫不等式（Chebyshev's inequality），分别说明非负性、阈值为正和二阶矩条件。
系统整理柯西--施瓦茨不等式（Cauchy--Schwarz inequality）及詹森不等式（Jensen's inequality）的概率形式与可积条件。
切尔诺夫界（Chernoff bound）和霍夫丁不等式（Hoeffding's inequality）作为选读，分别检查指数矩与独立有界条件。

\section{随机变量序列的收敛方式}
\label{sec:ps-convergence-modes}
定义\emph{几乎必然收敛}（almost sure convergence）、\emph{依概率收敛}（convergence in probability）、
\emph{依分布收敛}（convergence in distribution）及\emph{均方收敛}（mean-square convergence）。
建立常用蕴含关系并配反例；依分布收敛在极限分布函数的连续点定义。
说明依分布收敛到常数可推出依概率收敛，但依概率收敛本身不保证期望收敛。

\section{弱大数定律与强大数定律}
\label{sec:ps-laws-large-numbers}
对独立同分布随机变量 $X_1,X_2,\ldots$，记 $\overline X_n=n^{-1}\sum_{j=1}^nX_j$。
先在共同方差有限时用切比雪夫不等式证明\emph{弱大数定律}（weak law of large numbers）。
再陈述独立同分布且 $\mathbb E[|X_1|]<\infty$ 时的\emph{强大数定律}（strong law of large numbers），其完整证明列为选读。
区分依概率与几乎必然两种结论，不把一次样本轨迹的单调改善视为定理内容。

\section{中心极限定理与正态近似}
\label{sec:ps-central-limit}
对独立同分布、均值为 $\mu$、方差为 $\sigma^2\in(0,\infty)$ 的随机变量，陈述\emph{中心极限定理}（central limit theorem）：
\begin{equation}
  \frac{\sqrt n(\overline X_n-\mu)}{\sigma}
  \xrightarrow{\mathrm d} N(0,1),
  \label{eq:ps-clt}
\end{equation}
其中 $\xrightarrow{\mathrm d}$ 表示依分布收敛，$N(0,1)$ 表示标准正态分布。
用特征函数给出选读证明，讨论二项分布的正态近似与连续性修正。
强调近似质量与原分布及目标尾部有关，不设一个对所有分布都可靠的固定样本量阈值。

\section{连续映射、斯卢茨基定理与德尔塔方法}
\label{sec:ps-asymptotic-tools}
建立\emph{连续映射定理}（continuous mapping theorem）与\emph{斯卢茨基定理}（Slutsky's theorem），明确函数连续性与另一个因子依概率趋于常数等条件。
在渐近正态与可微条件下介绍\emph{德尔塔方法}（delta method），非退化的一阶近似另要求相关导数非零。
由此解释用一致估计的标准差替代总体标准差的依据，为统计推断作准备。

\section{其他近似与适用范围}
\label{sec:ps-approximation-boundaries}
讨论二项分布在试验次数趋于无穷、成功概率趋于零且其乘积趋于有限正数时的泊松极限。
通过有限总体抽样、依赖序列与无限方差分布，说明不能直接套用独立同分布的有限方差中心极限定理。
比较有限样本概率界与渐近近似：前者给可核查的界，后者描述适当极限下的分布形状。

\section{例题、习题与学习检查}
规划用切比雪夫不等式控制样本均值误差、用中心极限定理近似计数概率、构造依概率收敛但期望不收敛的例子。
证明题训练各收敛方式的关系；应用题要求逐条写明独立性、同分布和矩条件。
学习检查以区分“大数定律描述趋近目标”“中心极限定理描述标准化波动”为主。





